\chapter{Spin, Measurement, and the Two-Dimensional State Space}

\lectureoverview{A spin is directional, yet every measurement of one of its components returns only \(+1\) or \(-1\). We first formulate that operational puzzle, then see why ordinary set logic fails for incompatible measurements. A two-dimensional complex vector space supplies the right language: probabilities come from amplitudes, distinguishable alternatives are orthogonal, and the up--down, right--left, and in--out bases encode three spatial directions without turning a state vector into an ordinary arrow.}

\section{From Qubit to Spin}

The two-state system introduced as a qubit will now be called a \emph{spin}. The change of name brings spatial direction into the discussion, but it also creates a danger in the word \emph{vector}. Two different objects must be kept apart:

\begin{definition}
A \emph{three-vector} is an ordinary real vector in three-dimensional space. An unqualified \emph{vector} will mean an element of an abstract vector space, such as the complex state space used in quantum mechanics.
\end{definition}

A spin can be prepared relative to an axis in ordinary space. Its state, however, will be represented by an abstract vector. The relation between those two statements is one of the things we must discover; it cannot be assumed from a picture.

\section{Why the Apparatus Cannot Be Removed}

Classical mechanics encourages us to imagine an ideal observer who learns about a system while disturbing it as little as desired. Suppose, for example, that we want to see a macroscopic object. We illuminate it, collect the reflected light with a lens, and amplify the image. The light exerts some force, but in the classical idealization its intensity can be reduced while the amplification is increased. The disturbance can be made arbitrarily small without abandoning the idea of a measurement.

That limiting procedure fails in quantum mechanics. An apparatus can be improved, shielded, and calibrated, but a measurement cannot in general be made both informative and arbitrarily noninvasive. Something about the system changes. This does not mean that every measurement destroys every property or that the outcome cannot be reproduced. It means that the apparatus is part of the physical question being asked, and the state after the measurement cannot be ignored.

For the moment, let the apparatus be a black box with two features: it can be oriented in space, and its display records either \(+1\) or \(-1\). A small arrow on the board is a convenient label for a prepared spin:

\begin{figure}[htbp]
\centering
\includegraphics[width=0.64\textwidth]{lecture_02_figure_02.png}
\caption{The first provisional mark for a spin prepared along an axis. The arrow labels a preparation; it is not yet a classical model of the state.}
\lectureframe{00:05:15}
\end{figure}

Choose the axes so that \(z\) points up the board, \(x\) points to the right, and \(y\) points out of the board. If the apparatus is oriented along a unit direction \(\hat m\), its result is
\begin{equation}
\sigma_{\hat m}=\pm1.
\end{equation}
In particular,
\[
\sigma_z=\pm1,\qquad \sigma_x=\pm1,\qquad \sigma_y=\pm1.
\]
At this stage, the word \emph{component} is operational: \(\sigma_z\) is what the oriented apparatus reports when it points along \(z\). It is not yet the \(z\)-component of a hidden classical arrow.

\subsection{Repeatability and disturbance}

Suppose a \(z\)-measurement gives \(+1\). With no intervening operation, a second \(z\)-measurement gives \(+1\) again, as do further repetitions. The first measurement has prepared a state in which that same question has a definite answer.

\begin{classroomqa}
\audiencequestion{If a quantum measurement disturbs the spin, why does an immediate repetition return the same result?}
\lecturerresponse{The disturbance is not indiscriminate. Measuring \(\sigma_z\) can change what would happen in a later \(x\)- or \(y\)-measurement while leaving the newly measured value of \(\sigma_z\) reproducible. A useful experiment requires exactly this limited form of stability.}
\end{classroomqa}

Now turn the apparatus through \(180^\circ\). A spin that gave \(+1\) with the upward orientation gives \(-1\) with the inverted orientation. Turn the apparatus upright again and the result is \(+1\). Orientation therefore matters. A thermometer would not report minus the temperature merely because it was turned upside down; the spin analyzer behaves differently because it asks a directional question.

\begin{figure}[htbp]
\centering
\includegraphics[width=0.72\textwidth]{lecture_02_figure_03.png}
\caption{The board records the two-valued labels \(\sigma_z=\pm1\) and \(\sigma_x=\pm1\), together with the \(x\)-axis cue.}
\lectureframe{00:12:46}
\end{figure}

The contrast with a classical unit vector is immediate. Classical components vary continuously between \(-1\) and \(+1\). A unit vector perpendicular to the \(x\)-axis has \(x\)-component zero. A spin analyzer pointed along \(x\), by contrast, never displays zero. It displays \(+1\) or \(-1\), even when the spin was prepared along \(z\).

\section{Single Outcomes and Vector-Like Averages}

Prepare a spin by obtaining \(\sigma_z=+1\), and then measure \(\sigma_x\). Half the identically prepared trials give \(+1\), and half give \(-1\). Thus
\begin{equation}
\langle \sigma_x\rangle=0.
\end{equation}
The same preparation gives
\begin{equation}
\langle \sigma_y\rangle=0.
\end{equation}
Conversely, a spin prepared with \(\sigma_x=+1\) has equal probabilities for the two outcomes of a subsequent \(z\)- or \(y\)-measurement.

There are two distinct experiments hidden in that statement. Repeatedly measuring \(\sigma_x\) on the \emph{same} spin after the first \(x\)-result merely confirms that result. To obtain the fifty--fifty statistics, one needs many spins prepared in the same \(z\)-state, or one spin repeatedly returned to that same preparation before each \(x\)-measurement.

\begin{classroomqa}
\audiencequestion{What happens if a spin has not been deliberately prepared before it enters the apparatus?}
\lecturerresponse{The result is still \(+1\) or \(-1\), but its probabilities depend on whatever state the spin actually has. Preparation is the controlled procedure that makes repeated trials comparable. If an unknown manipulation occurs between two measurements, repeatability is no longer guaranteed; if nothing occurs, the same-component result is reproducible.}
\end{classroomqa}

\subsection{Two arbitrary directions}

Prepare the spin along a unit direction \(\hat n\) by obtaining \(+1\) from an analyzer oriented along \(\hat n\). Then measure along another unit direction \(\hat m\). Each trial still gives only \(\pm1\), but the average over identically prepared trials is
\begin{equation}
\boxed{\langle\sigma_{\hat m}\rangle=\hat n\cdot\hat m
=\cos\theta_{nm}.}
\label{eq:spin-average}
\end{equation}
The blackboard uses the shorter notation \(\langle\sigma_m\rangle=n\cdot m\):

\begin{figure}[htbp]
\centering
\includegraphics[width=0.74\textwidth]{lecture_02_figure_04.png}
\caption{The vector-like law belongs to the ensemble average, not to an individual \(\pm1\) outcome.}
\lectureframe{00:23:22}
\end{figure}

Equation~\eqref{eq:spin-average} includes the earlier cases. If \(\hat m=\hat n\), the average is \(+1\); since no result can exceed \(+1\), every trial gives \(+1\). If \(\hat m=-\hat n\), every trial gives \(-1\). Perpendicular axes give zero average and equal frequencies.

If \(p_+\) and \(p_-\) denote the probabilities for the two possible results, then
\[
p_++p_-=1,\qquad p_+-p_-=\hat n\cdot\hat m.
\]
Consequently,
\begin{equation}
p_\pm=\frac{1\pm\hat n\cdot\hat m}{2}.
\label{eq:direction-probabilities}
\end{equation}
\editorialnote{The lecture states the two outcomes and their mean. Equation~\eqref{eq:direction-probabilities} is the direct algebraic consequence of those two facts.}

\begin{classroomqa}
\audiencequestion{Is the quantity in angle brackets an \emph{expected value} or an ordinary average?}
\lecturerresponse{It is the mean of the recorded results over many repetitions. \emph{Expectation value} is standard terminology, but no extra kind of averaging is intended here. The apparatus produces a list of \(+1\)'s and \(-1\)'s; the brackets denote the arithmetic mean of that list in the large-ensemble limit.}
\end{classroomqa}

\begin{classroomqa}
\audiencequestion{Did the arrow \(\hat n\) already have a geometrical meaning before the average law was known?}
\lecturerresponse{Initially it was a label for how the preparation apparatus was oriented. The inversion experiment and the dependence of the average on \(\hat n\cdot\hat m\) give that label operational geometry. The apparatus defines what the components mean; a single result never supplies a classical projection.}
\end{classroomqa}

This is the first reconciliation of the directional picture with the discrete data. The individual spin is not behaving as a classical arrow with continuously valued components. Nevertheless, the statistics of many identically prepared spins transform exactly like the components of a unit three-vector.

\section{Classical Propositions as Sets}

Before constructing the quantum state space, consider how logic works for a classical system. Let the possible states form a set \(\Omega\). A proposition \(A\) is represented by the subset of states for which \(A\) is true.

For a six-sided die,
\[
A=\{2,4,6\}
\]
may represent \emph{the result is even}, while
\[
B=\{4,5,6\}
\]
may represent \emph{the result is greater than three}. Logical operations become set operations:
\begin{equation}
\neg A=A^c,\qquad
A\land B=A\cap B,\qquad
A\lor B=A\cup B.
\label{eq:classical-logic}
\end{equation}
Here \(\lor\) is the inclusive OR: \(A\lor B\) is true when \(A\) is true, \(B\) is true, or both are true. It is not the exclusive everyday use of \emph{either \(A\) or \(B\), but not both}.

\begin{classroomqa}
\audiencequestion{A spin measurement has only two outcomes, but the preparation direction varies continuously. Is the state set finite or continuous?}
\lecturerresponse{The set of possible preparation directions is continuous even though every chosen analyzer has only two readouts. This is precisely why counting detector labels is not enough to describe the space of states.}
\end{classroomqa}

For classical subsets, intersection and union do not care which proposition is mentioned first:
\[
A\cap B=B\cap A,\qquad A\cup B=B\cup A.
\]
An experimental test of a classical compound proposition can therefore be idealized without building an order into the logic.

\section{An Order-Dependent OR Test}

Consider the quantum propositions
\begin{equation}
A:\sigma_z=+1,\qquad B:\sigma_x=+1.
\end{equation}
Each proposition has a clear test. Orient the apparatus along \(z\) to test \(A\), or along \(x\) to test \(B\). The trouble appears when one tries to verify \(A\lor B\) by testing one proposition and, if it is false, testing the other.

Imagine that an unseen preparer has supplied a \(z\)-up spin. If \(A\) is tested first, the result is certainly \(+1\), and the OR test ends with \emph{true}.

Reverse the order. Testing \(B\) first gives
\[
P(B)=\frac12,\qquad P(\neg B)=\frac12.
\]
If \(B\) is false, the \(x\)-measurement has nevertheless prepared a definite \(x\)-state. Testing \(A\) on that state gives
\[
P(A\mid \neg B)=\frac12,\qquad
P(\neg A\mid \neg B)=\frac12.
\]
The reversed procedure therefore reports both propositions false with probability
\begin{equation}
P(\neg B\ \text{then}\ \neg A)
=\frac12\cdot\frac12=\frac14.
\end{equation}
One ordering always accepts the OR; the other rejects it in one quarter of the trials.

\begin{classroomqa}
\audiencequestion{Is the contradiction caused merely by our incomplete knowledge of the secretly prepared state?}
\lecturerresponse{No. The hidden preparation was chosen only to make the comparison sharp. The first \(x\)-measurement physically prepares an \(x\)-state, so the later \(z\)-test is being made on a changed state. Giving the observer more information about the initial preparation does not undo that intervention.}
\end{classroomqa}

\begin{classroomqa}
\audiencequestion{Could the OR be tested by a more elaborate hierarchy of questions rather than by this sequential rule?}
\lecturerresponse{Different experimental procedures can certainly be defined, but then they are different quantum measurements. There is no procedure-independent subset \(A\cup B\) whose truth is simply uncovered without regard to the apparatus. The order dependence is the warning that classical set logic is not the native logic of these propositions.}
\end{classroomqa}

Later, propositions will be represented by projection operators. The concise algebraic warning is
\begin{equation}
P_A P_B\ne P_B P_A.
\end{equation}
Nothing more from operator theory is needed yet. What matters is that incompatible tests do not commute, and their compound propositions cannot be treated as ordinary set unions.

\begin{classroomqa}
\audiencequestion{Can a two-particle version of this experiment be arranged so that the choice of test acts backward in time?}
\lecturerresponse{That question requires machinery not yet developed here. No conclusion about backward causation follows from the one-spin sequence, so the issue is left open rather than guessed at.}
\end{classroomqa}

\section{The Linear State Space}

A set contains elements, but there is generally no natural way to add two elements and obtain a third. A linear vector space has more structure. If \(\ket{A}\) and \(\ket{B}\) are vectors and \(\alpha,\beta\in\mathbb C\), then
\begin{equation}
\alpha\ket{A}+\beta\ket{B}
\end{equation}
is another vector in the same space. Quantum states are represented in such a complex linear space.

For each ket \(\ket{A}\) there is a corresponding dual bra \(\bra{A}\). Pairing a bra with a ket gives a complex number, the inner product. It obeys conjugate symmetry,
\begin{equation}
\braket{A}{B}=\braket{B}{A}^{*}.
\label{eq:conjugate-symmetry}
\end{equation}

\begin{figure}[htbp]
\centering
\includegraphics[width=0.78\textwidth]{lecture_02_figure_05.png}
\caption{The board summarizes linear combinations and the conjugate reversal of an inner product.}
\lectureframe{00:46:53}
\end{figure}

The squared length of a vector is
\begin{equation}
\|\ket{A}\|^2=\braket{A}{A}.
\end{equation}
Equation~\eqref{eq:conjugate-symmetry} makes this number equal to its own complex conjugate, hence real; the axioms of the inner product also make it nonnegative. Two vectors are orthogonal when
\begin{equation}
\braket{A}{B}=0.
\end{equation}
The dimension of the space is the largest number of mutually orthogonal nonzero directions that can be chosen.

\begin{classroomqa}
\audiencequestion{Must every basis consist of mutually orthogonal vectors?}
\lecturerresponse{No. A basis only has to be linearly independent and span the space. Orthogonal bases are especially convenient because coefficients and probabilities are easy to extract. In two dimensions, two nonparallel vectors may form a basis even when their inner product is not zero.}
\end{classroomqa}

\section{The Two-Dimensional Spin Space}

The state space of one spin is postulated to be a two-dimensional complex vector space,
\begin{equation}
\mathcal H_{\mathrm{spin}}\cong\mathbb C^2.
\end{equation}
There are continuously many possible state vectors, but only two independent basis vectors are required to express any of them.

Begin with the states prepared by a \(z\)-oriented apparatus:
\begin{align}
\ket{u}&\longleftrightarrow \sigma_z=+1,\\
\ket{d}&\longleftrightarrow \sigma_z=-1.
\end{align}
The letters \(u,d\), the signs \(+,-\), and the arrows \(\uparrow,\downarrow\) are alternative labels for the same two preparations. Since a \(z\)-measurement distinguishes them with certainty, their state vectors are orthogonal:
\begin{equation}
\braket{u}{d}=0.
\end{equation}
Choose them to have unit length,
\[
\braket{u}{u}=\braket{d}{d}=1.
\]

Every spin state can then be written
\begin{equation}
\ket{A}=\alpha_u\ket{u}+\alpha_d\ket{d},
\qquad \alpha_u,\alpha_d\in\mathbb C.
\label{eq:general-spin-state}
\end{equation}
The probability rule is
\begin{equation}
P(u|A)=|\alpha_u|^2,\qquad
P(d|A)=|\alpha_d|^2.
\label{eq:born-z}
\end{equation}
Because one of the two outcomes must occur,
\begin{equation}
|\alpha_u|^2+|\alpha_d|^2=1.
\label{eq:normalization}
\end{equation}

\begin{figure}[htbp]
\centering
\includegraphics[width=0.78\textwidth]{lecture_02_figure_06.png}
\caption{A general up--down expansion, the sum of the two outcome probabilities, and the equivalent unit-norm condition.}
\lectureframe{01:03:38}
\end{figure}

The equivalence with unit normalization follows directly:
\begin{align}
\braket{A}{A}
&=(\alpha_u^*\bra{u}+\alpha_d^*\bra{d})
  (\alpha_u\ket{u}+\alpha_d\ket{d})\\
&=|\alpha_u|^2\braket{u}{u}
 +\alpha_u^*\alpha_d\braket{u}{d}\\
&\quad +\alpha_d^*\alpha_u\braket{d}{u}
 +|\alpha_d|^2\braket{d}{d}\\
&=|\alpha_u|^2+|\alpha_d|^2=1.
\end{align}
Thus a physical pure state is represented by a normalized vector.

\begin{classroomqa}
\audiencequestion{Are these statements being derived, or are they postulates?}
\lecturerresponse{At this point they are the rules being fitted to the experiments. The two-dimensional space, the probability rule, and the interpretation of orthogonality are being introduced heuristically. Their mutual consistency can be checked now; the systematic relation between states and observables comes next.}
\end{classroomqa}

\subsection{Orthogonality as distinguishability}

The physical content of orthogonality is more useful than the picture of perpendicular arrows:

\begin{quote}
States that can be distinguished with certainty by one measurement are represented by orthogonal vectors.
\end{quote}

Up and down satisfy this criterion. If the state is up, the \(z\)-apparatus never reports down, and vice versa. A nonzero overlap indicates that one state can sometimes produce an outcome associated with the other.

\section{The Right--Left Basis}

Let \(\ket{r}\) and \(\ket{l}\) denote the states with definite \(x\)-spin:
\[
\ket{r}\longleftrightarrow\sigma_x=+1,\qquad
\ket{l}\longleftrightarrow\sigma_x=-1.
\]
Because the space is two-dimensional, each must be a linear combination of \(\ket{u}\) and \(\ket{d}\). A right-prepared spin gives up and down with equal probability, so the two coefficients must have equal magnitude. Choose the simplest phase convention:
\begin{equation}
\ket{r}=\frac{1}{\sqrt2}\bigl(\ket{u}+\ket{d}\bigr).
\label{eq:right-state}
\end{equation}
Right and left are distinguishable with certainty by the \(x\)-apparatus, so the left state must be orthogonal to the right state. One convenient choice is
\begin{equation}
\ket{l}=\frac{1}{\sqrt2}\bigl(\ket{u}-\ket{d}\bigr).
\label{eq:left-state}
\end{equation}
Indeed,
\begin{align}
\braket{l}{r}
&=\frac12(\bra{u}-\bra{d})(\ket{u}+\ket{d})\\
&=\frac12(1-1)=0.
\end{align}

\begin{figure}[htbp]
\centering
\includegraphics[width=0.82\textwidth]{lecture_02_figure_07.png}
\caption{The completed right--left construction: equal up--down magnitudes, opposite relative signs, and zero inner product.}
\lectureframe{01:12:15}
\end{figure}

The overall minus sign of \(\ket{l}\) is conventional. Replacing \(\ket{l}\) by \(-\ket{l}\) changes none of its probabilities. Adding and subtracting Equations~\eqref{eq:right-state} and~\eqref{eq:left-state} gives the inverse change of basis:
\begin{align}
\ket{u}&=\frac{1}{\sqrt2}\bigl(\ket{r}+\ket{l}\bigr),\\
\ket{d}&=\frac{1}{\sqrt2}\bigl(\ket{r}-\ket{l}\bigr).
\end{align}
The symmetry is important. Starting from the \(x\)-basis, an up or down state is an equal-weight superposition of right and left, just as a right or left state is an equal-weight superposition of up and down.

\section{The In--Out Basis and Relative Phase}

Orient the apparatus along \(y\), perpendicular to the board. Write
\[
\ket{i}\longleftrightarrow\sigma_y=+1,\qquad
\ket{o}\longleftrightarrow\sigma_y=-1,
\]
where the ket label \(i\) means \emph{in} and the unadorned \(i\) multiplying a ket remains the imaginary unit.

Equal probabilities for up and down require equal coefficient magnitudes. Orthogonality between in and out requires a relative phase that makes their overlap vanish. A symmetric choice is
\begin{align}
\ket{i}&=\frac{1}{\sqrt2}\bigl(\ket{u}+i\ket{d}\bigr),\\
\ket{o}&=\frac{1}{\sqrt2}\bigl(\ket{u}-i\ket{d}\bigr).
\label{eq:in-out-states}
\end{align}
In the \((u,d)\) basis, all six states are
\begin{align}
\ket{u}&=\binom10,
&
\ket{d}&=\binom01,
\\
\ket{r}&=\frac1{\sqrt2}\binom11,
&
\ket{l}&=\frac1{\sqrt2}\binom1{-1},
\\
\ket{i}&=\frac1{\sqrt2}\binom1{i},
&
\ket{o}&=\frac1{\sqrt2}\binom1{-i}.
\end{align}

\begin{figure}[htbp]
\centering
\includegraphics[width=0.82\textwidth]{lecture_02_figure_08.png}
\caption{The six column vectors in the up--down basis. The lower pair differs by the sign of the relative imaginary phase.}
\lectureframe{01:20:50}
\end{figure}

Complex conjugation in the bra is essential to the orthogonality check:
\begin{align}
\braket{i}{o}
&=\frac12(\bra{u}-i\bra{d})(\ket{u}-i\ket{d})\\
&=\frac12\bigl(1+(-i)(-i)\bigr)=0.
\end{align}

\begin{classroomqa}
\audiencequestion{Does choosing \(+i\) for \emph{in} rather than \emph{out} encode the handedness of the spatial axes?}
\lecturerresponse{A convention has entered, and its geometrical meaning will become clearer when the Pauli matrices and rotations are developed. For now, the pair must be orthogonal and symmetrically related to the other bases; exchanging the labels or reversing the \(y\)-axis changes the convention, not the consistency of the construction.}
\end{classroomqa}

\subsection{Checking symmetry in the right--left basis}

Use
\[
\ket{u}=\frac{\ket{r}+\ket{l}}{\sqrt2},\qquad
\ket{d}=\frac{\ket{r}-\ket{l}}{\sqrt2}
\]
in the in-state:
\begin{align}
\ket{i}
&=\frac1{\sqrt2}
\left(
\frac{\ket{r}+\ket{l}}{\sqrt2}
+i\frac{\ket{r}-\ket{l}}{\sqrt2}
\right)\\
&=\frac{1+i}{2}\ket{r}
 +\frac{1-i}{2}\ket{l}.
\label{eq:in-in-rl}
\end{align}
Although neither coefficient visibly contains \(1/\sqrt2\),
\begin{equation}
\left|\frac{1+i}{2}\right|^2
=\left|\frac{1-i}{2}\right|^2
=\frac12.
\end{equation}

\begin{figure}[htbp]
\centering
\includegraphics[width=0.82\textwidth]{lecture_02_figure_09.png}
\caption{The live change of basis for the in-state and the check that each right--left probability is \(1/2\).}
\lectureframe{01:29:10}
\end{figure}

Thus each state in one of the three bases gives equal probabilities when measured in either of the other two bases. The \(z\), \(x\), and \(y\) bases are mutually unbiased.

\section{Why There Are Three Distinguished Basis Pairs}

\begin{classroomqa}
\audiencequestion{Can there be a fourth pair of states, as symmetrically related to these three pairs as they are to one another?}
\lecturerresponse{Not for a single two-state system. Once the three mutually unbiased pairs have been fixed, there is no fourth pair with the same equal-overlap relation. The fact that there are three such spatial labels is tied to the three dimensions of ordinary space, even though the complex state space itself has dimension two.}
\end{classroomqa}

This connection can be seen by counting physical parameters. A general column
\[
\binom{\alpha_u}{\alpha_d}
\]
contains two complex numbers, hence four real numbers. Normalization removes one:
\[
|\alpha_u|^2+|\alpha_d|^2=1.
\]
Multiplying both entries by a common phase changes no measurement probability:
\begin{equation}
\ket{A}\sim e^{i\chi}\ket{A}.
\label{eq:global-phase}
\end{equation}
That removes one more real parameter. Therefore
\begin{equation}
4-1-1=2
\end{equation}
physical real parameters remain, exactly the number of angles needed to specify a direction in three-dimensional space.

A convenient representative makes the relation explicit:
\begin{equation}
\ket{\psi(\theta,\phi)}
=\cos\frac{\theta}{2}\,\ket{u}
+e^{i\phi}\sin\frac{\theta}{2}\,\ket{d}.
\label{eq:bloch-parameterization}
\end{equation}
\editorialnote{Equation~\eqref{eq:bloch-parameterization} is the standard normalized representative of the two-parameter family identified by the lecture. The explicit half-angle form is included to make the parameter count concrete; its dynamical and rotational interpretation belongs to later lectures.}

\begin{classroomqa}
\audiencequestion{What is the precise mathematics of ignoring the overall phase?}
\lecturerresponse{All normalized vectors on the circle \(\{e^{i\chi}\ket{\psi}\}\) represent the same physical state. One identifies those vectors as equivalent, much as points related by a declared redundancy are treated as one physical point. Relative phase between components remains measurable; only a phase multiplying the entire vector is discarded.}
\end{classroomqa}

\begin{classroomqa}
\audiencequestion{May non-normalized vectors or vectors with negative components still appear in a calculation?}
\lecturerresponse{Yes. Linear algebra allows them, and an intermediate vector need not yet represent a prepared physical state. A nonzero vector can be normalized at the end. Negative and complex components are amplitudes, not negative probabilities; probabilities arise from squared magnitudes after the state has been normalized.}
\end{classroomqa}

\begin{classroomqa}
\audiencequestion{What becomes of this picture in higher-dimensional spaces, or for more complicated spinors?}
\lecturerresponse{The same linear principles continue, but the relation between state-space dimension, spatial rotations, and observable spin becomes richer. That extension is outside the present two-state construction.}
\end{classroomqa}

\section{Orthogonality Is Not Spatial Perpendicularity}

The word \emph{orthogonal} now carries an abstract, operational meaning. It must not be confused with the angle between two arrows drawn in ordinary space. The spatial directions up and down are opposite, yet the state vectors \(\ket{u}\) and \(\ket{d}\) are orthogonal. The spatial directions up and right are perpendicular, yet
\begin{equation}
\braket{u}{r}=\frac1{\sqrt2}\ne0.
\end{equation}

\begin{classroomqa}
\audiencequestion{How can opposite spatial directions be orthogonal while perpendicular spatial directions are not?}
\lecturerresponse{State-space overlap measures experimental distinguishability, not the Euclidean angle between the labels. Up and down can be told apart with certainty by one \(z\)-measurement, so their overlap vanishes. Up and right cannot: a right-prepared spin sometimes gives the up result. Their state vectors therefore overlap.}
\end{classroomqa}

The same idea has a classical probabilistic analogue. Two probability distributions with disjoint support are perfectly distinguishable: an observed result identifies which distribution was used. Distributions with overlapping support cannot be distinguished without error from one sample. Quantum orthogonality plays the role of nonoverlap, while a nonzero inner product signals that some outcomes are compatible with either preparation.

\begin{classroomqa}
\audiencequestion{Should the two-dimensional state space be visualized as the physical plane containing the up and right arrows?}
\lecturerresponse{No. The labels \(x,y,z\) refer to orientations of the apparatus in ordinary space. The kets belong to a complex vector space with its own inner product. The parameter count connects the two structures, but identifying their pictures point by point would confuse spatial perpendicularity with state-space orthogonality.}
\end{classroomqa}

\section{States, Measurements, and What Remains to Be Done}

The construction has so far proceeded by fitting vectors to operational facts. The three apparatus orientations label three pairs of normalized states, and the overlaps reproduce the observed probabilities. A rigorous treatment must still say how an observable such as \(\sigma_z\) acts on a state and why its possible outcomes are its eigenvalues. That is the next step.

The earlier logical puzzle can now be stated more sharply. Propositions correspond to subspaces or their projectors, and incompatible projectors do not commute. The order-sensitive OR test was not a verbal paradox; it was an early appearance of this noncommutative structure.

\begin{classroomqa}
\audiencequestion{If the state is not an ordinary arrow, why do the three labels \(x,y,z\) and the average \(\langle\sigma_{\hat m}\rangle\) look so classical?}
\lecturerresponse{The classical vector appears in ensemble averages. A single spin gives a discrete answer, while many identically prepared spins yield means that transform like components of a direction. Macroscopic classical behavior emerges when large collections make those averages stable; it should not be read backward as a hidden continuous value in each individual spin.}
\end{classroomqa}

\begin{classroomqa}
\audiencequestion{Is the common phase redundancy an example of gauge invariance?}
\lecturerresponse{It is a simple invariance under a common phase, and it foreshadows broader gauge ideas. Developing that connection here would outrun the present construction, so only the physically relevant statement is retained: common phase is invisible, relative phase is not.}
\end{classroomqa}

Finally, a measurement should not be pictured as merely uncovering a value that the spin carried in the classical sense all along. It prepares the state associated with the recorded result. An immediate repetition of the same measurement is therefore consistent, while a new ensemble of identically prepared spins displays the probability distribution dictated by the relevant amplitudes. The mathematical description of that preparation and readout will become precise once observables are introduced as operators.
